\section{Chinese}

This is the ancient mythological layer of China, told with its gods and its first beings. It is older and more pictorial than the philosophical Taoism treated elsewhere in this collection. Here the cosmos hatches from an egg, a giant holds the sky and earth apart, and a serpent-bodied mother shapes the first people from clay. The philosophical reading of the same pattern belongs to the Taoism section.

\subsection{The Creation Model}

In the beginning there was only a formless \textbf{chaos}, shaped like an egg. All things were mingled together inside it, the light and the heavy, the clear and the murky, none of them yet apart. Within the egg a first being slept and grew for a vast age.

When the first being woke, it broke the egg open. The light and clear parts floated up and became the \textbf{sky}. The heavy and murky parts sank down and became the \textbf{earth}. The first being stood between them and held them apart. Each day the sky rose higher, the earth grew thicker, and the being grew taller to match, pushing the two ever further apart so they could not fall back together. This went on for another vast age.

When the work was done, the first being lay down and died. Its body became the world. Its breath became the wind and clouds. Its voice became the thunder. Its two eyes became the sun and the moon. Its limbs became the mountains, its blood the rivers, its flesh the soil, and its hair the stars and the plants.

The world was made but empty of people. A shaping mother walked the new earth. She took yellow clay and modeled the first humans by hand in her own image. When she tired of shaping each one by hand, she dipped a cord into the mud and flung it, and every drop that fell became a person. So the world was filled.

\subsection{The Model with Its Names}

The ancient Chinese accounts are gathered from later collections, chiefly the \textbf{Huainanzi} and the \textbf{Shanhaijing} (the Classic of Mountains and Seas), along with other Han-era texts. The myths do not form one scripture. They form a set of linked images.

\subsubsection{Hundun, the Primordial Chaos}

Before anything there was the formless primordial chaos (\textbf{Hundun}). It is the undivided state in which all opposites are mingled and none are yet distinct. It is pictured as a featureless mass or as the cosmic egg. Hundun is the starting condition that every later distinction comes out of.

\subsubsection{Pangu and the Cosmic Egg}

\textbf{Pangu} is the first being who sleeps inside the egg for eighteen thousand years while \textbf{yin} and \textbf{yang} lie mingled and undivided around him.

\begin{itemize}
 \item When Pangu wakes, he breaks the egg. The clear and light \textbf{yang} rises to form the heavens. The heavy and murky \textbf{yin} sinks to form the earth.
 \item Pangu stands between sky and earth and holds them apart. He grows taller every day for another eighteen thousand years, keeping the two from collapsing back into one.
 \item When the work is done, Pangu dies and his body becomes the substance of the world. His breath is the wind, his voice the thunder, his left eye the sun and his right eye the moon, his limbs the mountains, his blood the rivers, his flesh the soil, and his hair the stars.
\end{itemize}

So the splitting of the one into the polar two, and the holding of the two apart, is the whole work of creation. The world is the first body distributed into many forms.

\subsubsection{Nuwa, Shaper and Mender}

\textbf{Nuwa} is the serpent-bodied goddess who fills and repairs the world.

\begin{itemize}
 \item She shapes the first humans from \textbf{yellow clay}, modeling them by hand in her own image. When she tires, she flings a mud-soaked cord, and the falling drops become more people. One reading even sorts the careful hand-made from the cord-flung into the noble and the common.
 \item When the pillar of heaven is broken and the sky cracks, Nuwa \textbf{repairs the heavens}. She melts five-colored stones to patch the sky, and she cuts the legs from a great turtle to prop up the four corners of the world.
\end{itemize}

The shaper of life is also the mender of the broken order. Nuwa makes the people and keeps the sky standing over them.

\subsubsection{The Sovereigns and the Founders}

After the first making, the tradition tells of the \textbf{Three Sovereigns and Five Emperors}, half-divine culture-bringers who taught humanity the arts of civilization.

\begin{longtable}{L{0.28\linewidth}
L{0.6\linewidth}}
\toprule
\textbf{Figure} & \textbf{Gift to humanity} \\
\midrule
\textbf{Fuxi} & Hunting, fishing, writing, and the eight trigrams \\
\textbf{Nuwa} & The making and renewal of humankind \\
\textbf{Shennong} & Farming, the plough, and herbal medicine \\
\textbf{Huangdi} (the Yellow Emperor) & Governance, the calendar, and many crafts \\
\bottomrule
\end{longtable}

These figures bridge the world of the first beings and the world of human kings and history. The exact list of the Three Sovereigns varies between sources.

\subsection{A Note on the Philosophical Layer}

The same pattern of one source unfolding into a polar two, then a sustaining three, then the many, is stated in abstract form by the later Taoist tradition. The well-known line runs that the source gives birth to the One, the One to the Two, the Two to the Three, and the Three to the ten thousand things. The mythological Pangu story and the philosophical Tao account are two readings of one shape. For the philosophical development, see the \textbf{Taoism} section in this collection.

\subsection{Comparison}

The ancient Chinese model is the clearest single image of the move from undivided source to ordered world. A formless egg-chaos splits into the polar pair of sky and earth, a standing third holds them apart, and the one body becomes the many things. Unlike the Mesopotamian world, the world here is made not by combat but by separation and self-sacrifice, and humanity is shaped from clay not to serve the gods but to fill the empty earth.
